\section{Experimental Setup}\label{sec:setup}

\subsection{Data}\label{sec:data}
We train and test on the image set of EgoPoint-Ground~\cite{egopoint}, which we call EgoPoint, with its official split. Its paper reports that 39\% of the images are real captures with smart glasses and the rest are synthetically composed or edited; the release does not mark which is which, so, like other work on this benchmark~\cite{li2026vistaref,li2026pointvgr}, we report on the whole test split. After converting the labels (Appendix~\ref{app:data}), the test split has 2,604 images of the wearer pointing at a target and 425 images without a hand. We dropped 43 test images whose fingertip is not annotated, which changes the reported accuracy by less than one point. EgoPoint has no frames with a non-pointing hand and none with pointing but no target, so we evaluate the hand gate but not the pointing and target gates. In about 70\% of the images, the fingertip lies inside the target box.

\paragraph{COCO overlap.} The edited images were made by adding a pointing hand to photographs from COCO~\cite{lin2014coco}, selected through RefCOCO~\cite{yu2016modeling}, and the release names each source photograph. At least 41\% of our test images (1,242 of 3,029) are COCO train2014 photographs, which every COCO-pretrained detector has seen in training, together with box labels for the objects they contain. We therefore separate COCO-pretrained baselines from those that, like SpatialNet, never saw COCO data, and test whether the overlap inflates their scores (Section~\ref{sec:comparison}).

\subsection{Metrics}\label{sec:metrics}
The main metric is \emph{P@0.5}, the share of test images whose predicted target box overlaps the labelled box with IoU${\geq}0.5$, the standard grounding metric on this benchmark~\cite{egopoint,li2026vistaref,li2026pointvgr}; P@0.7 and the mean IoU (mIoU) measure how tight the boxes are. \emph{Centre in target}, the share of images whose predicted target centre lies inside the labelled box, measures whether the model picked the right object, independently of how well it drew the box. In a \emph{recognition test}, an off-the-shelf zero-shot recogniser chooses among the 348 target category names of EgoPoint from a crop of each model's predicted box, as a ``what is that?'' assistant would. We use two recognisers trained with different objectives, CLIP ViT-B/16~\cite{radford2021clip} and SigLIP base~\cite{zhai2023siglip}, without any training, and a crop counts as recognised if the top-ranked name is the labelled category (Appendix~\ref{app:recog}). Differences between models are tested with paired bootstrap intervals over the test images (1,000 resamples). For ablations, a difference counts as an effect only if its interval excludes zero \emph{and} it exceeds the change that retraining the full model with another seed produces.

\subsection{Baselines}\label{sec:baselines}
\paragraph{Published large models.} We copy the reported results of VistaRef~\cite{li2026vistaref}, PointVG-R and the multimodal LLMs that it evaluates~\cite{li2026pointvgr}. We did not re-run them, and their protocols differ slightly from ours (Appendix~\ref{app:comparison}), so we read differences from them as large gaps, not precise ones.

\paragraph{Compact detectors at the same compute.} YOLO26n~\cite{ultralytics2026yolo26} at a 128\,px input (0.225 GFLOPs) and NanoDet-Plus-m~\cite{rangilyu2021nanodet} at 160\,px (0.218 GFLOPs) cost about as much compute as SpatialNet (0.231 GFLOPs). Each is trained on the same split to output one box for the wearer's hand and one for the target, and is scored by the same code on the same images. Besides a benchmark of both detectors from several starting points (Appendix~\ref{app:comparison}), we trained YOLO26n four times under one fully specified recipe: from ImageNet weights, the pretraining of our own backbone, with two seeds; from ImageNet weights with our distractor pasting; and from COCO weights.

\subsection{Implementation and Deployment}\label{sec:device}
We train with AdamW and a one-cycle schedule for up to 180 epochs, with early stopping on the validation loss, on one NVIDIA T4 (1\,h\,40\,min, seed 42); Appendix~\ref{app:method} lists all settings. Retraining with seeds 1 and 2 gives the same target P@0.5 (58.9--59.0\% in all three runs), and the spread across seeds calibrates the single-run comparisons below.

We deploy the 160\,px model, exported to ONNX and FP16, on an engineering unit of a pair of smart glasses: a pre-production test device, not a retail product. It has a smart-glasses system-on-chip with a low-power microcontroller, a 12\,MP ultra-wide camera in portrait orientation, a wireless link to a companion phone and a 160\,mAh battery. The model squashes the portrait frames to its square input, as it does with the mostly landscape training images. Because the unit's power management is not tuned for production, we treat it as a measurement platform and compare two operating modes on the same unit rather than report battery life: \emph{on the glasses}, SpatialNet runs continuously on camera frames; in \emph{offloading}, the glasses only capture frames and send them to the phone at 0.5 frames per second, where localisation would run on the phone or in the cloud.
